\begin{table}[H]
\centering
\caption{Cod floor classes under alternative quantile conventions ($n=24$ fit residuals) and the implied constants.}
\label{tab:si-quantiles}
\small
\begin{tabular}{lrrrrr}
\toprule
convention & $e_{q05}$ & $e_{q10}$ & $C^*$ (q10) & $C_{\mathrm{vac}}$ (q10) & $C_{\mathrm{vac}}$ (q05) \\
\midrule
linear interpolation (declared) & -287.36 & -80.87 & 91.59 & 215.22 & 8.72 \\
lower order statistic & -323.51 & -82.56 & 89.90 & 213.53 & -27.42 \\
inverted CDF & -323.51 & -82.56 & 89.90 & 213.53 & -27.42 \\
Hazen & -325.15 & -106.65 & 65.81 & 189.43 & -29.06 \\
Weibull & -327.61 & -203.03 & -30.57 & 93.05 & -31.52 \\
\bottomrule
\end{tabular}
\par\smallskip\begin{minipage}{0.97\linewidth}\footnotesize Sorted fit residuals (kt): -329.0, -323.5, -82.6, -76.9, -26.7, -10.9, -7.6, -6.2, -5.6, -3.2, 0.2, 1.0, 2.6, 3.7, 4.2, 4.5, 4.7, 11.0, 16.2, 33.4, 88.8, 103.0, 131.2, 206.5. The two collapse residuals (destination years 1992 and 1993) are the only ones below $-172.5$ kt; the declared q10 floor sits between the second and third order statistics, which is why the convention matters for q10 and not for q05.\end{minipage}
\end{table}
